% =====================================================================
% H3: Early, small checkpoints become usable embeddings (owner: ...)
% Code and results: experiments/layerwise_probe/
% =====================================================================
\section{H3: Early, small checkpoints become usable embeddings}
\label{sec:h3}

\paragraph{Hypothesis.}
MIPIC \citep{mipic2026} adds a progressive information chaining (PIC) loss to Matryoshka representation learning (MRL; \citealp{kusupati2022mrl}).
PIC chains six checkpoints, written as (layer, width): $(1,16)$, $(2,32)$, $(3,64)$, $(4,256)$, $(5,512)$ and $(6,768)$.
At each checkpoint it links the first $d$ coordinates of the [CLS] vector at layer $l$ to the next checkpoint.
The paper therefore implies that the early, small checkpoints become useful embeddings in their own right, but it evaluates only the final layer.
We test the vector that PIC trains, the raw [CLS] prefix $h_l[0{:}d]$, not \texttt{pooler\_output}.
The claim is: at $(1,16)$ and $(2,32)$, MIPIC beats both MRL and untrained TinyBERT, and MRL+PIC beats MRL.
The last comparison isolates PIC, because MRL+PIC differs from MRL only by the PIC loss.

\paragraph{Data.}
The main task is the repository's 3-class tweet sentiment probe data (26,732 training and 3,432 test tweets).
These are the Kaggle Tweet Sentiment Extraction data, not the TweetEval benchmark.
As a secondary check, we use the STS-B test split.
We remove the 118 of its 1,379 pairs that appear in the training corpus, which leaves 1,261 clean pairs.

\paragraph{Models and probes.}
We compare the four trained variants (MRL, MRL+SIA, MRL+PIC, MIPIC), each trained with seeds 42, 1 and 2 and attention KL off, and use their \texttt{best.pt} checkpoints (Section~\ref{sec:setup}).
Untrained TinyBERT \citep{jiao-etal-2020-tinybert} is the reference, so a gain must be over the backbone and not only over MRL.
On frozen float32 features, we fit a logistic-regression probe with the same settings as the repository's evaluation script (no feature scaling, $C{=}1$, at most 1,000 iterations) at every layer and width.
We report the six PIC checkpoints.

\paragraph{Metrics.}
\begin{itemize}
  \item \textbf{Macro-F1} on the test tweets: the F1 score of each class, averaged with equal weight.
  \emph{Intuition:} every sentiment class counts the same, so a probe that mostly predicts the most common class scores low even if its accuracy looks fine.

  \item \textbf{Paired bootstrap difference.}
  For two models A and B, we redraw the 3,432 test tweets with replacement 1,000 times, use the same draws for both models, and compute macro-F1(A) $-$ macro-F1(B) on each draw.
  We report the observed difference and, as its 95\% CI, the 2.5th and 97.5th percentiles of the redrawn differences.
  \emph{Intuition:} the CI shows whether the gap would survive a different sample of test tweets of the same size; if it excludes zero, the gap is not luck in the test set.
  It does not cover seed variation, which the decision rule handles by requiring every seed.

  \item \textbf{STS-B Spearman} ($\times$100): the rank correlation between the gold similarity scores and the cosine similarity of the two sentences' prefixes.
  \emph{Intuition:} it asks whether the prefix ranks sentence pairs by meaning the way humans do, with no classifier trained on top.
\end{itemize}

\paragraph{Decision rules.}
We fixed the rule before we computed any results.
In one seed, H3 holds if MIPIC beats both MRL and untrained TinyBERT at $(1,16)$ and $(2,32)$ with 95\% CIs excluding zero, and MRL+PIC also beats MRL at both checkpoints (six comparisons).
Across seeds, H3 is \emph{supported} if this holds in all three seeds, \emph{partially supported} if it holds in exactly two, and \emph{not supported} otherwise.
STS-B is descriptive and does not enter the rule.

\begin{table}[t]
\centering
\small
\setlength{\tabcolsep}{4pt}
\resizebox{\columnwidth}{!}{%
\begin{tabular}{@{}lrrrrrc@{}}
\toprule
$(l,d)$ & Base & MRL & MRL+PIC & MIPIC & MIPIC $-$ MRL & CI excl.\ 0 \\
\midrule
$(1,16)$  & 33.88 & 33.97 & 44.62 & \textbf{46.19} & $+$4.69 to $+$20.43 & 3/3 \\
$(2,32)$  & 57.20 & 56.98 & 61.11 & \textbf{62.06} & $+$4.63 to $+$5.98 & 3/3 \\
$(3,64)$  & \textbf{68.24} & 67.50 & 67.98 & 67.33 & $-$0.73 to $+$0.53 & 0/3 \\
$(4,256)$ & \textbf{72.04} & 71.92 & 70.25 & 69.62 & $-$3.37 to $-$1.41 & 3/3 \\
$(5,512)$ & \textbf{73.07} & 72.26 & 71.02 & 69.90 & $-$3.08 to $-$1.46 & 3/3 \\
$(6,768)$ & \textbf{73.51} & 71.92 & 70.44 & 69.47 & $-$2.77 to $-$2.02 & 3/3 \\
\bottomrule
\end{tabular}%
}
\caption{H3: sentiment macro-F1 of a probe on the raw [CLS] prefix $h_l[0{:}d]$ at the six PIC checkpoints $(l,d)$. Model columns are means over seeds 42, 1 and 2; base is a single untrained model. MIPIC $-$ MRL is the paired bootstrap difference, as a range over the three seeds; the last column counts the seeds whose 95\% CI excludes zero. The best value in each row is in bold.}
\label{tab:h3}
\end{table}

\paragraph{Results.}
Table~\ref{tab:h3} shows the main results.
\textbf{H3 is supported}: all six required comparisons pass in 3 of 3 seeds.
At $(1,16)$, MIPIC is above MRL by $+$4.69 to $+$20.43 points and above untrained TinyBERT by $+$5.02 to $+$20.77; at $(2,32)$, by $+$4.63 to $+$5.98 and $+$4.19 to $+$5.88.
MRL+PIC beats MRL by $+$6.00 to $+$16.21 at $(1,16)$ and $+$2.48 to $+$6.45 at $(2,32)$; the narrowest case (MRL+PIC at $(2,32)$, seed 42) still has a lower CI bound of $+$0.73.
The gain is large at $(1,16)$, smaller at $(2,32)$ and gone at $(3,64)$, where no seed shows a significant difference.
From $(4,256)$ on it reverses: at every remaining checkpoint and in every seed, MIPIC is significantly below MRL (by 1.41 to 3.37 points) and below untrained TinyBERT (by 1.83 to 5.19 points).
MRL+SIA alone shows no early gain ($-$0.68 to $+$0.15 over MRL at $(1,16)$, with no CI excluding zero), so the effect comes from PIC.
The direction is stable, but the size varies a lot by seed: MIPIC minus MRL at $(1,16)$ is $+$4.69 in seed 1, $+$11.54 in seed 2 and $+$20.43 in seed 42.
With attention KL on (seed 42 only), the rule still passes: MIPIC is $+$8.84 [$+$6.98, $+$10.61] above MRL and $+$9.18 [$+$7.28, $+$10.91] above untrained TinyBERT at $(1,16)$, and $+$4.07 and $+$3.98 at $(2,32)$, all with CIs above zero; MRL+PIC has no attention KL term, so its comparisons do not change.
But the $(1,16)$ gain over MRL falls from $+$20.43 to $+$8.84.
The early sentiment gain does not carry over to STS-B.
At $(2,32)$, untrained TinyBERT is best in every seed (59.29, against 37.49 to 52.37 for MIPIC and 37.72 to 48.34 for MRL+PIC).
At $(1,16)$, MIPIC is above it only in seed 42 (39.97 against 25.19) and below it in seeds 1 and 2 (20.88 and 21.20).
The probe fit at $(5,512)$ reached the iteration limit for untrained TinyBERT, for MRL+SIA in all seeds and for MRL in seed 2, so that row partly reflects the solver.
As a side result, at the final layer untrained TinyBERT matches or beats every trained model on this sentiment probe, both on \texttt{pooler\_output} and on the raw [CLS] at $(6,768)$ in Table~\ref{tab:h3}; we return to this in Section~\ref{sec:discussion}.

\paragraph{Takeaway.}
PIC does what its design suggests at the earliest checkpoints: in every seed, the first 16 and 32 [CLS] coordinates of layers 1 and 2 carry more sentiment information than in MRL or the untrained model, and SIA alone does not do this.
The gain is paid for later in the network, where MIPIC falls below both MRL and the backbone from layer 4 on, and it does not extend to sentence similarity.
Its size depends strongly on the seed and shrinks with attention KL on, so it should be quoted as a range.
